\documentclass[10pt,a4paper]{amsart}
\usepackage{amsmath,amssymb,amsthm,graphicx}
\usepackage{verbatim}
\newenvironment{reference}{\comment}{\endcomment}
%\newenvironment{reference}{}{}
\newenvironment{slogan}{\comment}{\endcomment}
\newenvironment{history}{\comment}{\endcomment}
% For commutative diagrams we use Xy-pic
\usepackage[all]{xy}

% We use 2cell for 2-commutative diagrams.
\xyoption{2cell}
% Theorem environments.
%
\theoremstyle{plain}
\newtheorem{theorem}[subsection]{Theorem}
\newtheorem{proposition}[subsection]{Proposition}
\newtheorem{lemma}[subsection]{Lemma}

\theoremstyle{definition}
\newtheorem{definition}[subsection]{Definition}
\newtheorem{example}[subsection]{Example}
\newtheorem{exercise}[subsection]{Exercise}
\newtheorem{situation}[subsection]{Situation}

\theoremstyle{remark}
\newtheorem{remark}[subsection]{Remark}
\newtheorem{remarks}[subsection]{Remarks}

\numberwithin{equation}{subsection}

% Macros
%
\def\lim{\mathop{\mathrm{lim}}\nolimits}
\def\colim{\mathop{\mathrm{colim}}\nolimits}
\def\Spec{\mathop{\mathrm{Spec}}}
\def\Hom{\mathop{\mathrm{Hom}}\nolimits}
\def\Ext{\mathop{\mathrm{Ext}}\nolimits}
\def\SheafHom{\mathop{\mathcal{H}\!\mathit{om}}\nolimits}
\def\SheafExt{\mathop{\mathcal{E}\!\mathit{xt}}\nolimits}
\def\Sch{\mathit{Sch}}
\def\Mor{\mathop{\mathrm{Mor}}\nolimits}
\def\Ob{\mathop{\mathrm{Ob}}\nolimits}
\def\Sh{\mathop{\mathit{Sh}}\nolimits}
\def\NL{\mathop{N\!L}\nolimits}
\def\CH{\mathop{\mathrm{CH}}\nolimits}
\def\proetale{{pro\text{-}\acute{e}tale}}
\def\etale{{\acute{e}tale}}
\def\QCoh{\mathit{QCoh}}
\def\Ker{\mathop{\mathrm{Ker}}}
\def\Im{\mathop{\mathrm{Im}}}
\def\Coker{\mathop{\mathrm{Coker}}}
\def\Coim{\mathop{\mathrm{Coim}}}

% Boxtimes
%
\DeclareMathSymbol{\boxtimes}{\mathbin}{AMSa}{"02}

%
% Macros for moduli stacks/spaces
%
\def\QCohstack{\mathcal{QC}\!\mathit{oh}}
\def\Cohstack{\mathcal{C}\!\mathit{oh}}
\def\Spacesstack{\mathcal{S}\!\mathit{paces}}
\def\Quotfunctor{\mathrm{Quot}}
\def\Hilbfunctor{\mathrm{Hilb}}
\def\Curvesstack{\mathcal{C}\!\mathit{urves}}
\def\Polarizedstack{\mathcal{P}\!\mathit{olarized}}
\def\Complexesstack{\mathcal{C}\!\mathit{omplexes}}
% \Pic is the operator that assigns to X its picard group, usage \Pic(X)
% \Picardstack_{X/B} denotes the Picard stack of X over B
% \Picardfunctor_{X/B} denotes the Picard functor of X over B
\def\Pic{\mathop{\mathrm{Pic}}\nolimits}
\def\Picardstack{\mathcal{P}\!\mathit{ic}}
\def\Picardfunctor{\mathrm{Pic}}
\def\Deformationcategory{\mathcal{D}\!\mathit{ef}}
\usepackage{cleveref}
\crefname{originaldefinition}{original definition}{original definitions}
\Crefname{originaldefinition}{Original definition}{Original definitions}
\makeatletter
\newlabel{definition-bivariant-class}{{0B76}{source}{}{sourcepointer-1}{}}
\newlabel{definition-bivariant-class@cref}{{[originaldefinition][0][]0B76}{[1][0][]source}{}{}{}}
\makeatother
\crefname{originaldefinition}{original definition}{original definitions}
\Crefname{originaldefinition}{Original definition}{Original definitions}
\makeatletter
\newlabel{definition-gysin-homomorphism}{{02T8}{source}{}{sourcepointer-2}{}}
\newlabel{definition-gysin-homomorphism@cref}{{[originaldefinition][0][]02T8}{[1][0][]source}{}{}{}}
\makeatother
\crefname{originallemma}{original lemma}{original lemmas}
\Crefname{originallemma}{Original lemma}{Original lemmas}
\makeatletter
\newlabel{lemma-rational-equivalence-family}{{02S4}{source}{}{sourcepointer-3}{}}
\newlabel{lemma-rational-equivalence-family@cref}{{[originallemma][0][]02S4}{[1][0][]source}{}{}{}}
\makeatother
\crefname{originalremark}{original remark}{original remarks}
\Crefname{originalremark}{Original remark}{Original remarks}
\makeatletter
\newlabel{remark-gysin-on-cycles}{{0B6Z}{source}{}{sourcepointer-4}{}}
\newlabel{remark-gysin-on-cycles@cref}{{[originalremark][0][]0B6Z}{[1][0][]source}{}{}{}}
\makeatother
\crefname{originalreference}{original reference}{original references}
\Crefname{originalreference}{Original reference}{Original references}
\makeatletter
\newlabel{situation-setup}{{02QL}{source}{}{sourcepointer-5}{}}
\newlabel{situation-setup@cref}{{[originalreference][0][]02QL}{[1][0][]source}{}{}{}}
\makeatother
\title{A rule satisfying the bivariant condition for trivially normal bundles factors through rational equivalence}
\author{The Stacks Project Authors}
\date{stacks-project:chow.tex:local-snapshot-20261012; GNU FDL 1.2 or later}
\begin{document}
\maketitle
\noindent\textit{Editorial note.} All delivered mathematical text is the author's own, copied from the pinned Stacks Project chapter file chow.tex; no statement or proof step was rewritten, reordered or removed.\par\medskip
\noindent\textit{Editorial note.} Preamble provenance: the pinned chapter file opens with \textbackslash{}input\{preamble\} and carries no preamble text of its own, so the project's single Stacks-wide preamble staged at sources/stacks-chow/preamble.tex is excerpted verbatim as three spans, lines 8-12 (the verbatim package together with the project's own reference, slogan and history display-environment definitions, including the author's commented-out older form of the reference environment), lines 14-18 (Xy-pic setup) and lines 154-217 (the native theorem declarations and mathematical macros). The lines 8-12 span is required because a counted statement or its surrounding text may use those display environments, which the project typesets as comments. Nothing was added, removed or redefined; the wrapper class and the amsmath/amssymb/amsthm/graphicx base packages are supplied by the assembler.\par\medskip
\noindent\textit{Editorial note.} Same-chapter pointers are located at their \textbackslash{}label lines inside the pinned chapter file; cross-chapter pointers are located in the project's own tag list staged in the snapshot (sources/stacks-chow/tags.txt), at the verbatim lines carrying their canonical tags.\par\medskip
\noindent\textit{Editorial note.} The counted statement contains the author's own reference display environment at source line 6478 ('Very weak form of \textbackslash{}cite[Theorem 17.1]\{F\}'), which the project typesets as a comment and whose definition comes from the preamble span over lines 8-12 of the shared preamble. The citation is inside the commented block and never executed.\par\medskip
\noindent\textit{Allowed standard prerequisites (editorial).}\begin{itemize}
\item The bivariant class definition and the Gysin operation on the level of cycles. Same chapter chow.tex, labels definition-bivariant-class (proof line 6490), definition-gysin-homomorphism (proof line 6498) and remark-gysin-on-cycles (proof line 6501). The first two fix the setting, and the remark is what lets the operation c cap - be evaluated on individual cycles, which the statement needs in order to make sense.
\item The two fibre classes of a cycle on a product with the projective line. Same chapter chow.tex, label lemma-rational-equivalence-family, invoked at proof lines 6504 and 6515. Realises a rationally trivial cycle alpha as the difference of the two pullbacks of a single cycle beta on the product with the projective line, and later supplies the rational equivalence of the two pushes forward of c cap beta. That lemma is claimed separately as S14-algebra-rational-equiv-family.
\end{itemize}\medskip
\ifdefined\hypertarget\hypertarget{sourcepointer-1}{}\fi
\noindent\textit{Source reference.} Original definition 0B76: sources/stacks-chow/pinned-original-source.tex, source lines 6087--6087. Navigation pointer for \textbackslash{}ref\{definition-bivariant-class\} used inside the counted span (line 6490); the argument that consumes it is retained in the span.\par\medskip
\ifdefined\hypertarget\hypertarget{sourcepointer-2}{}\fi
\noindent\textit{Source reference.} Original definition 02T8: sources/stacks-chow/pinned-original-source.tex, source lines 5069--5069. Navigation pointer for \textbackslash{}ref\{definition-gysin-homomorphism\} used inside the counted span (line 6498); the argument that consumes it is retained in the span.\par\medskip
\ifdefined\hypertarget\hypertarget{sourcepointer-3}{}\fi
\noindent\textit{Source reference.} Original lemma 02S4: sources/stacks-chow/pinned-original-source.tex, source lines 3335--3335. Navigation pointer for \textbackslash{}ref\{lemma-rational-equivalence-family\} used inside the counted span (line 6504); the argument that consumes it is retained in the span.\par\medskip
\ifdefined\hypertarget\hypertarget{sourcepointer-4}{}\fi
\noindent\textit{Source reference.} Original remark 0B6Z: sources/stacks-chow/pinned-original-source.tex, source lines 5237--5237. Navigation pointer for \textbackslash{}ref\{remark-gysin-on-cycles\} used inside the counted span (line 6501); the argument that consumes it is retained in the span.\par\medskip
\ifdefined\hypertarget\hypertarget{sourcepointer-5}{}\fi
\noindent\textit{Source reference.} Original reference 02QL: sources/stacks-chow/pinned-original-source.tex, source lines 1372--1372. Navigation pointer for \textbackslash{}ref\{situation-setup\} used inside the counted span (line 6481); the argument that consumes it is retained in the span.\par\medskip
\begin{lemma}
\label{lemma-factors-through-rational-equivalence}
\begin{reference}
Very weak form of \cite[Theorem 17.1]{F}
\end{reference}
Let $(S, \delta)$ be as in Situation \ref{situation-setup}.
Let $f : X \to Y$ be a morphism of schemes locally of finite type over $S$.
Let $p \in \mathbf{Z}$. Suppose given a rule
which assigns to every locally of finite type morphism $Y' \to Y$
and every $k$ a map
$$
c \cap - : Z_k(Y') \longrightarrow \CH_{k - p}(X')
$$
where $Y' = X' \times_X Y$, satisfying condition (3) of
Definition \ref{definition-bivariant-class}
whenever $\mathcal{L}'|_{D'} \cong \mathcal{O}_{D'}$. Then
$c \cap -$ factors through rational equivalence.
\end{lemma}

\begin{proof}
The statement makes sense because given a triple
$(\mathcal{L}, s, i : D \to X)$ as in
Definition \ref{definition-gysin-homomorphism}
such that $\mathcal{L}|_D \cong \mathcal{O}_D$, then
the operation $i^*$ is defined on the level of cycles, see
Remark \ref{remark-gysin-on-cycles}.
Let $\alpha \in Z_k(X')$ be a cycle which is rationally equivalent to zero.
We have to show that $c \cap \alpha = 0$. By
Lemma \ref{lemma-rational-equivalence-family}
there exists a cycle $\beta \in Z_{k + 1}(X' \times \mathbf{P}^1)$
such that $\alpha = i_0^*\beta - i_\infty^*\beta$
where $i_0, i_\infty : X' \to X' \times \mathbf{P}^1$ are the
closed immersions of $X'$ over $0, \infty$. Since these are
examples of effective Cartier divisors with trivial normal
bundles, we see that $c \cap i_0^*\beta = j_0^*(c \cap \beta)$
and $c \cap i_\infty^*\beta = j_\infty^*(c \cap \beta)$
where $j_0, j_\infty : Y' \to Y' \times \mathbf{P}^1$ are
closed immersions as before. Since
$j_0^*(c \cap \beta) \sim_{rat} j_\infty^*(c \cap \beta)$
(follows from Lemma \ref{lemma-rational-equivalence-family}) we conclude.
\end{proof}

\end{document}
